\documentclass[10pt,a4paper]{amsart}
\usepackage{amsmath,amssymb,amsthm,graphicx}
\usepackage{verbatim}
\newenvironment{reference}{\comment}{\endcomment}
%\newenvironment{reference}{}{}
\newenvironment{slogan}{\comment}{\endcomment}
\newenvironment{history}{\comment}{\endcomment}

% For commutative diagrams we use Xy-pic
\usepackage[all]{xy}

% We use 2cell for 2-commutative diagrams.
\xyoption{2cell}
% Theorem environments.
%
\theoremstyle{plain}
\newtheorem{theorem}[subsection]{Theorem}
\newtheorem{proposition}[subsection]{Proposition}
\newtheorem{lemma}[subsection]{Lemma}

\theoremstyle{definition}
\newtheorem{definition}[subsection]{Definition}
\newtheorem{example}[subsection]{Example}
\newtheorem{exercise}[subsection]{Exercise}
\newtheorem{situation}[subsection]{Situation}

\theoremstyle{remark}
\newtheorem{remark}[subsection]{Remark}
\newtheorem{remarks}[subsection]{Remarks}

\numberwithin{equation}{subsection}

% Macros
%
\def\lim{\mathop{\mathrm{lim}}\nolimits}
\def\colim{\mathop{\mathrm{colim}}\nolimits}
\def\Spec{\mathop{\mathrm{Spec}}}
\def\Hom{\mathop{\mathrm{Hom}}\nolimits}
\def\Ext{\mathop{\mathrm{Ext}}\nolimits}
\def\SheafHom{\mathop{\mathcal{H}\!\mathit{om}}\nolimits}
\def\SheafExt{\mathop{\mathcal{E}\!\mathit{xt}}\nolimits}
\def\Sch{\mathit{Sch}}
\def\Mor{\mathop{\mathrm{Mor}}\nolimits}
\def\Ob{\mathop{\mathrm{Ob}}\nolimits}
\def\Sh{\mathop{\mathit{Sh}}\nolimits}
\def\NL{\mathop{N\!L}\nolimits}
\def\CH{\mathop{\mathrm{CH}}\nolimits}
\def\proetale{{pro\text{-}\acute{e}tale}}
\def\etale{{\acute{e}tale}}
\def\QCoh{\mathit{QCoh}}
\def\Ker{\mathop{\mathrm{Ker}}}
\def\Im{\mathop{\mathrm{Im}}}
\def\Coker{\mathop{\mathrm{Coker}}}
\def\Coim{\mathop{\mathrm{Coim}}}

% Boxtimes
%
\DeclareMathSymbol{\boxtimes}{\mathbin}{AMSa}{"02}

%
% Macros for moduli stacks/spaces
%
\def\QCohstack{\mathcal{QC}\!\mathit{oh}}
\def\Cohstack{\mathcal{C}\!\mathit{oh}}
\def\Spacesstack{\mathcal{S}\!\mathit{paces}}
\def\Quotfunctor{\mathrm{Quot}}
\def\Hilbfunctor{\mathrm{Hilb}}
\def\Curvesstack{\mathcal{C}\!\mathit{urves}}
\def\Polarizedstack{\mathcal{P}\!\mathit{olarized}}
\def\Complexesstack{\mathcal{C}\!\mathit{omplexes}}
% \Pic is the operator that assigns to X its picard group, usage \Pic(X)
% \Picardstack_{X/B} denotes the Picard stack of X over B
% \Picardfunctor_{X/B} denotes the Picard functor of X over B
\def\Pic{\mathop{\mathrm{Pic}}\nolimits}
\def\Picardstack{\mathcal{P}\!\mathit{ic}}
\def\Picardfunctor{\mathrm{Pic}}
\def\Deformationcategory{\mathcal{D}\!\mathit{ef}}
\usepackage{cleveref}
\crefname{originalreference}{original reference}{original references}
\Crefname{originalreference}{Original reference}{Original references}
\makeatletter
\newlabel{lemma-classes-points}{{0FHV}{source}{}{sourcepointer-1}{}}
\newlabel{lemma-classes-points@cref}{{[originalreference][0][]0FHV}{[1][0][]source}{}{}{}}
\makeatother
\crefname{originalreference}{original reference}{original references}
\Crefname{originalreference}{Original reference}{Original references}
\makeatletter
\newlabel{lemma-dim-0}{{0FHQ}{source}{}{sourcepointer-2}{}}
\newlabel{lemma-dim-0@cref}{{[originalreference][0][]0FHQ}{[1][0][]source}{}{}{}}
\makeatother
\crefname{originalreference}{original reference}{original references}
\Crefname{originalreference}{Original reference}{Original references}
\makeatletter
\newlabel{lemma-generated-by-separable}{{0FH5}{source}{}{sourcepointer-3}{}}
\newlabel{lemma-generated-by-separable@cref}{{[originalreference][0][]0FH5}{[1][0][]source}{}{}{}}
\makeatother
\crefname{originalreference}{original reference}{original references}
\Crefname{originalreference}{Original reference}{Original references}
\makeatletter
\newlabel{lemma-injective}{{09PV}{source}{}{sourcepointer-4}{}}
\newlabel{lemma-injective@cref}{{[originalreference][0][]09PV}{[1][0][]source}{}{}{}}
\makeatother
\crefname{originalreference}{original reference}{original references}
\Crefname{originalreference}{Original reference}{Original references}
\makeatletter
\newlabel{lemma-otimes}{{0FHX}{source}{}{sourcepointer-5}{}}
\newlabel{lemma-otimes@cref}{{[originalreference][0][]0FHX}{[1][0][]source}{}{}{}}
\makeatother
\crefname{originalreference}{original reference}{original references}
\Crefname{originalreference}{Original reference}{Original references}
\makeatletter
\newlabel{lemma-trace-disjoint-union}{{0FHP}{source}{}{sourcepointer-6}{}}
\newlabel{lemma-trace-disjoint-union@cref}{{[originalreference][0][]0FHP}{[1][0][]source}{}{}{}}
\makeatother
\crefname{originalreference}{original reference}{original references}
\Crefname{originalreference}{Original reference}{Original references}
\makeatletter
\newlabel{lemma-weil-additive}{{0FHJ}{source}{}{sourcepointer-7}{}}
\newlabel{lemma-weil-additive@cref}{{[originalreference][0][]0FHJ}{[1][0][]source}{}{}{}}
\makeatother
\crefname{originallemma}{original lemma}{original lemmas}
\Crefname{originallemma}{Original lemma}{Original lemmas}
\makeatletter
\newlabel{varieties-lemma-baby-stein}{{0FD1}{source}{}{sourcepointer-8}{}}
\newlabel{varieties-lemma-baby-stein@cref}{{[originallemma][0][]0FD1}{[1][0][]source}{}{}{}}
\makeatother
\crefname{originallemma}{original lemma}{original lemmas}
\Crefname{originallemma}{Original lemma}{Original lemmas}
\makeatletter
\newlabel{varieties-lemma-proper-geometrically-reduced-global-sections}{{0BUG}{source}{}{sourcepointer-9}{}}
\newlabel{varieties-lemma-proper-geometrically-reduced-global-sections@cref}{{[originallemma][0][]0BUG}{[1][0][]source}{}{}{}}
\makeatother
\title{Characterizations of the degree zero cohomology of a smooth projective scheme}
\author{The Stacks Project Authors}
\date{stacks-project:weil.tex:local-snapshot-20261012; GNU FDL 1.2 or later}
\begin{document}
\maketitle
\noindent\textit{Editorial note.} All delivered mathematical text is the author's own, copied from the pinned Stacks Project chapter file weil.tex; no statement or proof step was rewritten, reordered or removed.\par\medskip
\noindent\textit{Editorial note.} The counted proof is the author's complete argument: the reduction to an irreducible source, the equivalence of the first and third conditions through the axioms, and the finite separable base change making the chosen points rational.\par\medskip
\noindent\textit{Editorial note.} The label lemma-H-0-separable names the counted statement, and the proof environment immediately following it in the source is delivered verbatim as its proof.\par\medskip
\noindent\textit{Editorial note.} The counted statement is carried by the author's own lemma environment in the chapter source.\par\medskip
\noindent\textit{Editorial note.} Pre-open gate: python operations/agent-authoring/tools/precard\_gate.py --source weil.tex --label lemma-H-0-separable --start 3859 --end 3926 --local-id 28132 returned RESULT=OK (L28132 free, occupancy FREE, public interval clear). Pre-assembly tag self-check: canonical tag resolved and not present in reference/public-items.jsonl.\par\medskip
\noindent\textit{Editorial note.} The statement and proof are copied from the pinned chapter file; the cited inputs are kept as external pointers with their canonical Stacks tags.\par\medskip
\noindent\textit{Allowed standard prerequisites (editorial).}\begin{itemize}
\item The ring of global sections of a proper geometrically reduced scheme is a finite separable algebra. Varieties, the lemma on global sections of proper geometrically reduced schemes, and the lemma identifying the base change to the ring of global sections. Shows that the spectrum of the ring of global sections is smooth and projective, so that the statement makes sense, and allows reducing to a connected scheme.
\item Injectivity of the pullback to a closed point and equality of the classes of closed points. Same chapter, the lemma on injectivity of the map to the cohomology of a point and the lemma identifying the classes of points of the same degree. Provides the implication from the second condition to the first, and the equality of the pullbacks used in the reverse implication.
\item Behaviour of the degree zero part under base change of the field of global sections. Same chapter, the lemma comparing the cohomology of the base change with the tensor product over the degree zero cohomology of the base. Allows replacing the field of global sections by a finite separable extension so that the chosen points become rational.
\end{itemize}\medskip
\ifdefined\hypertarget\hypertarget{sourcepointer-1}{}\fi
\noindent\textit{Source reference.} Original reference 0FHV: sources/stacks-weil/pinned-original-source.tex, source lines 3697--3697. Prerequisite pointer invoked at proof line(s) 3918; it names a standard fact of the same chapter that the counted argument uses, and every step of the counted proof built on it is retained verbatim in the delivered body.\par\medskip
\ifdefined\hypertarget\hypertarget{sourcepointer-2}{}\fi
\noindent\textit{Source reference.} Original reference 0FHQ: sources/stacks-weil/pinned-original-source.tex, source lines 3432--3432. Prerequisite pointer invoked at proof line(s) 3924; it names a standard fact of the same chapter that the counted argument uses, and every step of the counted proof built on it is retained verbatim in the delivered body.\par\medskip
\ifdefined\hypertarget\hypertarget{sourcepointer-3}{}\fi
\noindent\textit{Source reference.} Original reference 0FH5: sources/stacks-weil/pinned-original-source.tex, source lines 2233--2233. Prerequisite pointer invoked at proof line(s) 3896; it names a standard fact of the same chapter that the counted argument uses, and every step of the counted proof built on it is retained verbatim in the delivered body.\par\medskip
\ifdefined\hypertarget\hypertarget{sourcepointer-4}{}\fi
\noindent\textit{Source reference.} Original reference 09PV: sources/stacks-weil/pinned-original-source.tex, source lines 3737--3737. Prerequisite pointer invoked at proof line(s) 3904; it names a standard fact of the same chapter that the counted argument uses, and every step of the counted proof built on it is retained verbatim in the delivered body.\par\medskip
\ifdefined\hypertarget\hypertarget{sourcepointer-5}{}\fi
\noindent\textit{Source reference.} Original reference 0FHX: sources/stacks-weil/pinned-original-source.tex, source lines 3753--3753. Prerequisite pointer invoked at proof line(s) 3909; it names a standard fact of the same chapter that the counted argument uses, and every step of the counted proof built on it is retained verbatim in the delivered body.\par\medskip
\ifdefined\hypertarget\hypertarget{sourcepointer-6}{}\fi
\noindent\textit{Source reference.} Original reference 0FHP: sources/stacks-weil/pinned-original-source.tex, source lines 3408--3408. Prerequisite pointer invoked at proof line(s) 3885; it names a standard fact of the same chapter that the counted argument uses, and every step of the counted proof built on it is retained verbatim in the delivered body.\par\medskip
\ifdefined\hypertarget\hypertarget{sourcepointer-7}{}\fi
\noindent\textit{Source reference.} Original reference 0FHJ: sources/stacks-weil/pinned-original-source.tex, source lines 2932--2932. Prerequisite pointer invoked at proof line(s) 3885; it names a standard fact of the same chapter that the counted argument uses, and every step of the counted proof built on it is retained verbatim in the delivered body.\par\medskip
\ifdefined\hypertarget\hypertarget{sourcepointer-8}{}\fi
\noindent\textit{Source reference.} Original lemma 0FD1: sources/stacks-chow/tags.txt, source lines 19201--19201. Prerequisite pointer invoked at proof line(s) 3893; it names a standard fact of a sibling chapter that the counted argument uses, and every step of the counted proof built on it is retained verbatim in the delivered body.\par\medskip
\ifdefined\hypertarget\hypertarget{sourcepointer-9}{}\fi
\noindent\textit{Source reference.} Original lemma 0BUG: sources/stacks-chow/tags.txt, source lines 14874--14874. Prerequisite pointer invoked at proof line(s) 3882; it names a standard fact of a sibling chapter that the counted argument uses, and every step of the counted proof built on it is retained verbatim in the delivered body.\par\medskip
\begin{lemma}
\label{lemma-H-0-separable}
Assume given (D0), (D1), (D2), and (D3) satisfying (A), (B), and (C).
Let $X$ be a smooth projective scheme over $k$.
Set $k' = \Gamma(X, \mathcal{O}_X)$. The following are equivalent
\begin{enumerate}
\item there exist finitely many closed points $x_1, \ldots, x_r \in X$
whose residue fields are separable over $k$ such that
$H^0(X) \to H^0(x_1) \oplus \ldots \oplus H^0(x_r)$ is injective,
\item the map $H^0(\Spec(k')) \to H^0(X)$ is an isomorphism.
\end{enumerate}
If this is true, then $H^0(X)$ is a finite separable algebra over $F$.
If $X$ is equidimensional of dimension $d$, then (1) and (2)
are also equivalent to
\begin{enumerate}
\item[(3)] the classes of closed points generate $H^{2d}(X)(d)$
as a module over $H^0(X)$.
\end{enumerate}
\end{lemma}

\begin{proof}
We observe that the statement makes sense because $k'$ is a finite separable
algebra over $k$ (Varieties, Lemma
\ref{varieties-lemma-proper-geometrically-reduced-global-sections})
and hence $\Spec(k')$ is smooth and projective over $k$.
The compatibility of $H^*$ with direct sums
(Lemmas \ref{lemma-weil-additive} and \ref{lemma-trace-disjoint-union})
shows that it suffices to prove the lemma when $X$ is connected.
Hence we may assume $X$ is irreducible and we have to show the
equivalence of (1), (2), and (3). Set $d = \dim(X)$.
This implies that $k'$ is a field finite separable
over $k$ and that $X$ is geometrically irreducible over $k'$, see
Varieties, Lemmas
\ref{varieties-lemma-proper-geometrically-reduced-global-sections} and
\ref{varieties-lemma-baby-stein}.

\medskip\noindent
By Lemma \ref{lemma-generated-by-separable} we see that the closed
points in (3) may be assumed to have separable residue fields over $k$.
By axioms (A)(a) and (A)(b) we see that conditions (1) and (3) are equivalent.

\medskip\noindent
If (2) holds, then pick any closed point $x \in X$ whose residue field
is finite separable over $k'$. Then
$H^0(\Spec(k')) = H^0(X) \to H^0(x)$ is injective for example by
Lemma \ref{lemma-injective}.

\medskip\noindent
Assume the equivalent conditions (1) and (3) hold. Choose
$x_1, \ldots, x_r \in X$ as in (1). Choose a finite separable
extension $k''/k'$. By Lemma \ref{lemma-otimes} we have
$$
H^0(X) \otimes_{H^0(\Spec(k'))} H^0(\Spec(k'')) =
H^0(X \times_{\Spec(k')} \Spec(k''))
$$
Thus in order to show that
$H^0(\Spec(k')) \to H^0(X)$ is an isomorphism
we may replace $k'$ by $k''$. Thus we may assume $x_1, \ldots, x_r$
are $k'$-rational points (this replaces each $x_i$ with multiple
points, so $r$ is increased in this step). By Lemma \ref{lemma-classes-points}
$\gamma(x_1) = \gamma(x_2) = \ldots = \gamma(x_r)$.
By axiom (A)(b) all the maps $H^0(X) \to H^0(x_i)$
are the same. This means (2) holds.

\medskip\noindent
Finally, Lemma \ref{lemma-dim-0} implies
$H^0(X)$ is a separable $F$-algebra if (1) holds.
\end{proof}

\end{document}
